% Original HER data-flow diagram. No relabeling edge creates real interactions.
\begin{tikzpicture}[x=\linewidth,y=1mm,
 font=\figurefont\fontsize{8.5}{11}\selectfont,
 box/.style={draw=BookRule,fill=BookPaper,line width=.8pt,
   text width=23mm,minimum height=11mm,align=center,inner sep=1.5mm},
 flow/.style={-{Latex[length=2mm]},draw=BookInk,line width=.9pt}]
 \node[box] (task) at (.14,0) {任务生成\\目标$g$};
 \node[box] (env) at (.50,0) {真实交互\\$\mu_g,P$};
 \node[box] (trajectory) at (.86,0) {轨迹\\$(s_t,a_t,s_{t+1})$};
 \node[box,fill=BookTeal!8] (relabel) at (.50,-26) {目标重标记\\$g'=(3,1)$};
 \node[box] (samples) at (.14,-26) {训练样本\\重算$r_t,\delta_t$};
 \draw[flow] (task) -- node[above] {采样} (env);
 \draw[flow] (env) -- node[above] {记录} (trajectory);
 \draw[flow] (trajectory.south) |- node[pos=.7,above] {已达位置} (relabel.east);
 \draw[flow] (relabel) -- (samples);
 \node[anchor=north,align=center,text=BookMuted] at (.5,-38)
   {上行：取得新经验\qquad 下行：重新解释已有经验};
\end{tikzpicture}
